%=====================================
%   rel-t-morse.tex
%   Truncated Microlocal Morse Theory
%   toshi2019
%=====================================

% -----------------------
% preamble
% -----------------------
% ここから本文 (\begin{document}) までの
% ソースコードに変更を加えた場合は
% 編集者まで連絡してください. 
% Don't change preamble code yourself. 
% If you add something
% (usepackage, newtheorem, newcommand, renewcommand),
% please tell it 
% to the editor of institutional paper of RUMS.

% ------------------------
% documentclass
% ------------------------
\documentclass[leqno, 12pt, a4paper, dvipdfmx]{article}

% ------------------------
% usepackage
% ------------------------
\usepackage{algorithm}
\usepackage{algorithmic}
\usepackage{amscd}
\usepackage{amsfonts}
\usepackage{amsmath}
\usepackage[psamsfonts]{amssymb}
\usepackage{amsthm}
\usepackage{ascmac}
\usepackage{bm}
\usepackage{color}
\usepackage{enumerate}
\usepackage{fancybox}
\usepackage[stable]{footmisc}
\usepackage{graphicx}
\usepackage{listings}
\usepackage{mathrsfs}
\usepackage{mathtools}
\usepackage{otf}
\usepackage{pifont}
\usepackage{proof}
\usepackage{subfigure}
\usepackage{tikz}
\usepackage{verbatim}
\usepackage[all]{xy}

\usetikzlibrary{cd}
\usetikzlibrary{arrows.meta}


% ================================
% パッケージを追加する場合のスペース 
\usepackage[dvipdfmx]{hyperref}
%\usepackage{xcolor}
\definecolor{darkgreen}{rgb}{0,0.45,0} 
\definecolor{darkred}{rgb}{0.75,0,0}
\definecolor{darkblue}{rgb}{0,0,0.6} 
\hypersetup{
    colorlinks=true,
    %citecolor=darkgreen,
    %linkcolor=darkred,
    %urlcolor=darkblue,
    citecolor=black,
    linkcolor=black,
    urlcolor=black,
}
\usepackage{pxjahyper}

%\usepackage[mathscr]{euscript} % mathscr のフォントを変更

%\usepackage{mmasym}

%=================================


% --------------------------
% theoremstyle
% --------------------------
\theoremstyle{definition}

% --------------------------
% newtheoem
% --------------------------

% 日本語で定理, 命題, 証明などを番号付きで用いるためのコマンドです. 
% If you want to use theorem environment in Japanece, 
% you can use these code. 
% Attention!
% All theorem enivironment numbers depend on 
% only section numbers.
\newtheorem{Axiom}{公理}%[section]
\newtheorem{Definition}[Axiom]{定義}
\newtheorem{Theorem}[Axiom]{定理}
\newtheorem{Proposition}[Axiom]{命題}
\newtheorem{Lemma}[Axiom]{補題}
\newtheorem{Corollary}[Axiom]{系}
\newtheorem{Example}[Axiom]{例}
\newtheorem{Claim}[Axiom]{主張}
\newtheorem{Property}[Axiom]{性質}
\newtheorem{Attention}[Axiom]{注意}
\newtheorem{Question}[Axiom]{問}
\newtheorem{Problem}[Axiom]{問題}
\newtheorem{Consideration}[Axiom]{考察}
\newtheorem{Alert}[Axiom]{警告}
\newtheorem{Fact}[Axiom]{事実}


% 日本語で定理, 命題, 証明などを番号なしで用いるためのコマンドです. 
% If you want to use theorem environment with no number in Japanese, You can use these code.
\newtheorem*{Axiom*}{公理}
\newtheorem*{Definition*}{定義}
\newtheorem*{Theorem*}{定理}
\newtheorem*{Proposition*}{命題}
\newtheorem*{Lemma*}{補題}
\newtheorem*{Example*}{例}
\newtheorem*{Corollary*}{系}
\newtheorem*{Claim*}{主張}
\newtheorem*{Property*}{性質}
\newtheorem*{Attention*}{注意}
\newtheorem*{Question*}{問}
\newtheorem*{Problem*}{問題}
\newtheorem*{Consideration*}{考察}
\newtheorem*{Alert*}{警告}
\newtheorem{Fact*}{事実}


% 英語で定理, 命題, 証明などを番号付きで用いるためのコマンドです. 
%　If you want to use theorem environment in English, You can use these code.
%all theorem enivironment number depend on only section number.
\newtheorem{Axiom+}{Axiom}%[section]
\newtheorem{Definition+}[Axiom+]{Definition}
\newtheorem{Theorem+}[Axiom+]{Theorem}
\newtheorem{Proposition+}[Axiom+]{Proposition}
\newtheorem{Lemma+}[Axiom+]{Lemma}
\newtheorem{Example+}[Axiom+]{Example}
\newtheorem{Corollary+}[Axiom+]{Corollary}
\newtheorem{Claim+}[Axiom+]{Claim}
\newtheorem{Property+}[Axiom+]{Property}
\newtheorem{Attention+}[Axiom+]{Attention}
\newtheorem{Question+}[Axiom+]{Question}
\newtheorem{Problem+}[Axiom+]{Problem}
\newtheorem{Consideration+}[Axiom+]{Consideration}
\newtheorem{Alert+}{Alert}
\newtheorem{Fact+}[Axiom+]{Fact}
\newtheorem{Remark+}[Axiom+]{Remark}

% ----------------------------
% commmand
% ----------------------------
% 執筆に便利なコマンド集です. 
% コマンドを追加する場合は下のスペースへ. 

% 集合の記号 (黒板文字)
\newcommand{\NN}{\mathbb{N}}
\newcommand{\ZZ}{\mathbb{Z}}
\newcommand{\QQ}{\mathbb{Q}}
\newcommand{\RR}{\mathbb{R}}
\newcommand{\CC}{\mathbb{C}}
\newcommand{\PP}{\mathbb{P}}
\newcommand{\KK}{\mathbb{K}}


% 集合の記号 (太文字)
\newcommand{\nn}{\mathbf{N}}
\newcommand{\zz}{\mathbf{Z}}
\newcommand{\qq}{\mathbf{Q}}
\newcommand{\rr}{\mathbf{R}}
\newcommand{\cc}{\mathbf{C}}
\newcommand{\pp}{\mathbf{P}}
\newcommand{\kk}{\mathbf{k}}

% 特殊な写像の記号
\newcommand{\ev}{\mathop{\mathrm{ev}}\nolimits} % 値写像
\newcommand{\pr}{\mathop{\mathrm{pr}}\nolimits} % 射影

% スクリプト体にするコマンド
%   例えば　{\mcal C} のように用いる
\newcommand{\mcal}{\mathcal}

% 花文字にするコマンド 
%   例えば　{\h C} のように用いる
\newcommand{\h}{\mathscr}

% ヒルベルト空間などの記号
\newcommand{\F}{\mcal{F}}
\newcommand{\X}{\mcal{X}}
\newcommand{\Y}{\mcal{Y}}
\newcommand{\Hil}{\mcal{H}}
\newcommand{\RKHS}{\Hil_{k}}
\newcommand{\Loss}{\mcal{L}_{D}}
\newcommand{\MLsp}{(\X, \Y, D, \Hil, \Loss)}

% 偏微分作用素の記号
\newcommand{\p}{\partial}

% 角カッコの記号 (内積は下にマクロがあります)
\newcommand{\lan}{\langle}
\newcommand{\ran}{\rangle}



% 圏の記号など
\newcommand{\Set}{{\bf Set}}
\newcommand{\Vect}{{\bf Vect}}
\newcommand{\FDVect}{{\bf FDVect}}
%\newcommand{\Ring}{{\bf Ring}}
\newcommand{\Ab}{{\bf Ab}}
\newcommand{\Mod}{\mathop{\mathrm{Mod}}\nolimits}
\newcommand{\Modf}{\mathop{\mathrm{Mod}^\mathrm{f}}\nolimits}
\newcommand{\CGA}{{\bf CGA}}
\newcommand{\GVect}{{\bf GVect}}
\newcommand{\Lie}{{\bf Lie}}
\newcommand{\dLie}{{\bf Liec}}



% 射の集合など
\newcommand{\Map}{\mathop{\mathrm{Map}}\nolimits} % 写像の集合
\newcommand{\Hom}{\mathop{\mathrm{Hom}}\nolimits} % 射集合
\newcommand{\End}{\mathop{\mathrm{End}}\nolimits} % 自己準同型の集合
\newcommand{\Aut}{\mathop{\mathrm{Aut}}\nolimits} % 自己同型の集合
\newcommand{\Mor}{\mathop{\mathrm{Mor}}\nolimits} % 射集合
\newcommand{\Ker}{\mathop{\mathrm{Ker}}\nolimits} % 核
\newcommand{\Img}{\mathop{\mathrm{Im}}\nolimits} % 像
\newcommand{\Cok}{\mathop{\mathrm{Coker}}\nolimits} % 余核
\newcommand{\Cim}{\mathop{\mathrm{Coim}}\nolimits} % 余像

% その他便利なコマンド
\newcommand{\dip}{\displaystyle} % 本文中で数式モード
\newcommand{\e}{\varepsilon} % イプシロン
\newcommand{\dl}{\delta} % デルタ
\newcommand{\pphi}{\varphi} % ファイ
\newcommand{\ti}{\tilde} % チルダ
\newcommand{\pal}{\parallel} % 平行
\newcommand{\op}{{\rm op}} % 双対を取る記号
\newcommand{\lcm}{\mathop{\mathrm{lcm}}\nolimits} % 最小公倍数の記号
\newcommand{\Probsp}{(\Omega, \F, \P)} 
\newcommand{\argmax}{\mathop{\rm arg~max}\limits}
\newcommand{\argmin}{\mathop{\rm arg~min}\limits}





% ================================
% コマンドを追加する場合のスペース 
%\newcommand{\OO}{\mcal{O}}



\makeatletter
\renewenvironment{proof}[1][\proofname]{\par
  \pushQED{\qed}%
  \normalfont \topsep6\p@\@plus6\p@\relax
  \trivlist
  \item[\hskip\labelsep
%        \itshape
         \bfseries
%    #1\@addpunct{.}]\ignorespaces
    {#1}]\ignorespaces
}{%
  \popQED\endtrivlist\@endpefalse
}
\makeatother

\renewcommand{\proofname}{Proof.}



%\renewcommand\proofname{\bf 証明} % 証明
\numberwithin{equation}{section}
\newcommand{\cTop}{\textsf{Top}}
%\newcommand{\cOpen}{\textsf{Open}}
\newcommand{\Op}{\mathop{\textsf{Op}}\nolimits}
\newcommand{\Ob}{\mathop{\textrm{Ob}}\nolimits}
\newcommand{\id}{\mathop{\mathrm{id}}\nolimits}
\newcommand{\pt}{\mathop{\mathrm{pt}}\nolimits}
\newcommand{\res}{\mathop{\rho}\nolimits}
\newcommand{\A}{\mcal{A}}
\newcommand{\B}{\mcal{B}}
\newcommand{\C}{\mcal{C}}
\newcommand{\D}{\mcal{D}}
\newcommand{\E}{\mcal{E}}
\newcommand{\G}{\mcal{G}}
%\newcommand{\H}{\mcal{H}}
\newcommand{\I}{\mcal{I}}
\newcommand{\J}{\mcal{J}}
\newcommand{\OO}{\mcal{O}}
\newcommand{\Ring}{\mathop{\textsf{Ring}}\nolimits}
\newcommand{\cAb}{\mathop{\textsf{Ab}}\nolimits}
%\newcommand{\Ker}{\mathop{\mathrm{Ker}}\nolimits}
\newcommand{\im}{\mathop{\mathrm{Im}}\nolimits}
\newcommand{\Coker}{\mathop{\mathrm{Coker}}\nolimits}
\newcommand{\Coim}{\mathop{\mathrm{Coim}}\nolimits}
\newcommand{\rank}{\mathop{\mathrm{rank}}\nolimits}
\newcommand{\Ht}{\mathop{\mathrm{Ht}}\nolimits}
\newcommand{\supp}{\mathop{\mathrm{supp}}\nolimits}
\newcommand{\colim}{\mathop{\mathrm{colim}}}
\newcommand{\Tor}{\mathop{\mathrm{Tor}}\nolimits}

\newcommand{\cat}{\mathscr{C}}

%筆記体
\newcommand{\cA}{\mcal{A}}
\newcommand{\cB}{\mcal{B}}
\newcommand{\cC}{\mcal{C}}
\newcommand{\cD}{\mcal{D}}
\newcommand{\cE}{\mcal{E}}
\newcommand{\cF}{\mcal{F}}
\newcommand{\cG}{\mcal{G}}
\newcommand{\cH}{\mcal{H}}
\newcommand{\cI}{\mcal{I}}
\newcommand{\cJ}{\mcal{J}}
\newcommand{\cK}{\mcal{K}}
\newcommand{\cL}{\mcal{L}}
\newcommand{\cM}{\mcal{M}}
\newcommand{\cN}{\mcal{N}}
\newcommand{\cO}{\mcal{O}}
\newcommand{\cP}{\mcal{P}}
\newcommand{\cQ}{\mcal{Q}}
\newcommand{\cR}{\mcal{R}}
\newcommand{\cS}{\mcal{S}}
\newcommand{\cT}{\mcal{T}}
\newcommand{\cU}{\mcal{U}}
\newcommand{\cV}{\mcal{V}}
\newcommand{\cW}{\mcal{W}}
\newcommand{\cX}{\mcal{X}}
\newcommand{\cY}{\mcal{Y}}
\newcommand{\cZ}{\mcal{Z}}


\newcommand{\scA}{\mathscr{A}}
\newcommand{\scB}{\mathscr{B}}
\newcommand{\scC}{\mathscr{C}}
\newcommand{\scD}{\mathscr{D}}
\newcommand{\scE}{\mathscr{E}}
\newcommand{\scF}{\mathscr{F}}
\newcommand{\scN}{\mathscr{N}}
\newcommand{\scO}{\mathscr{O}}
\newcommand{\scV}{\mathscr{V}}
\newcommand{\scU}{\mathscr{U}}


\newcommand{\ibA}{\mathop{\text{\textit{\textbf{A}}}}}
\newcommand{\ibB}{\mathop{\text{\textit{\textbf{B}}}}}
\newcommand{\ibC}{\mathop{\text{\textit{\textbf{C}}}}}
\newcommand{\ibD}{\mathop{\text{\textit{\textbf{D}}}}}
\newcommand{\ibE}{\mathop{\text{\textit{\textbf{E}}}}}
\newcommand{\ibF}{\mathop{\text{\textit{\textbf{F}}}}}
\newcommand{\ibG}{\mathop{\text{\textit{\textbf{G}}}}}
\newcommand{\ibH}{\mathop{\text{\textit{\textbf{H}}}}}
\newcommand{\ibI}{\mathop{\text{\textit{\textbf{I}}}}}
\newcommand{\ibJ}{\mathop{\text{\textit{\textbf{J}}}}}
\newcommand{\ibK}{\mathop{\text{\textit{\textbf{K}}}}}
\newcommand{\ibL}{\mathop{\text{\textit{\textbf{L}}}}}
\newcommand{\ibM}{\mathop{\text{\textit{\textbf{M}}}}}
\newcommand{\ibN}{\mathop{\text{\textit{\textbf{N}}}}}
\newcommand{\ibO}{\mathop{\text{\textit{\textbf{O}}}}}
\newcommand{\ibP}{\mathop{\text{\textit{\textbf{P}}}}}
\newcommand{\ibQ}{\mathop{\text{\textit{\textbf{Q}}}}}
\newcommand{\ibR}{\mathop{\text{\textit{\textbf{R}}}}}
\newcommand{\ibS}{\mathop{\text{\textit{\textbf{S}}}}}
\newcommand{\ibT}{\mathop{\text{\textit{\textbf{T}}}}}
\newcommand{\ibU}{\mathop{\text{\textit{\textbf{U}}}}}
\newcommand{\ibV}{\mathop{\text{\textit{\textbf{V}}}}}
\newcommand{\ibW}{\mathop{\text{\textit{\textbf{W}}}}}
\newcommand{\ibX}{\mathop{\text{\textit{\textbf{X}}}}}
\newcommand{\ibY}{\mathop{\text{\textit{\textbf{Y}}}}}
\newcommand{\ibZ}{\mathop{\text{\textit{\textbf{Z}}}}}

\newcommand{\ibx}{\mathop{\text{\textit{\textbf{x}}}}}

%\newcommand{\Comp}{\mathop{\mathrm{C}}\nolimits}
%\newcommand{\Komp}{\mathop{\mathrm{K}}\nolimits}
%\newcommand{\Domp}{\mathop{\mathsf{D}}\nolimits}%複体のホモトピー圏
%\newcommand{\Comp}{\mathrm{C}}
%\newcommand{\Komp}{\mathrm{K}}
%\newcommand{\Domp}{\mathsf{D}}%複体のホモトピー圏

\newcommand{\Comp}{\mathop{\mathrm{C}}\nolimits}
\newcommand{\Komp}{\mathop{\mathsf{K}}\nolimits}
\newcommand{\Domp}{\mathord{\mathsf{D}}}%\nolimits}
\newcommand{\Kompl}{\mathop{\mathsf{K}^\mathrm{+}}\nolimits}
\newcommand{\Kompu}{\mathop{\mathsf{K}^\mathrm{-}}\nolimits}
\newcommand{\Kompb}{\mathop{\mathsf{K}^\mathrm{b}}\nolimits}
\newcommand{\Dompl}{\mathop{\mathsf{D}^\mathrm{+}}\nolimits}
\newcommand{\Dompu}{\mathop{\mathsf{D}^\mathrm{-}}\nolimits}
\newcommand{\Dompb}{\mathop{\mathsf{D}^\mathrm{b}}\nolimits}
\newcommand{\Dompbf}{\mathop{\mathsf{D}_\mathrm{f}^\mathrm{b}}\nolimits}
\newcommand{\Domplb}{\mathop{\mathsf{D}^\mathrm{lb}}\nolimits}




\newcommand{\CCat}{\Comp(\cat)}
\newcommand{\KCat}{\Komp(\cat)}
\newcommand{\DCat}{\Domp(\cat)}%圏Cの複体のホモトピー圏
\newcommand{\HOM}{\mathop{\mathscr{H}\hspace{-2pt}om}\nolimits}%内部Hom
\newcommand{\RHOM}{\mathop{\mathrm{R}\hspace{-1.5pt}\HOM}\nolimits}

%\newcommand{\muS}{\mathop{\mathrm{SS}}\nolimits}
\newcommand{\muS}[1][]{\mathop{\mathrm{SS}}\nolimits_{#1}}
\newcommand{\RG}{\mathop{\mathrm{R}\hspace{-0pt}\Gamma}\nolimits}
\newcommand{\RHom}{\mathop{\mathrm{R}\hspace{-1.5pt}\Hom}\nolimits}
%\newcommand{\Rder}{\mathrm{R}}
\newcommand{\Rder}[1][]{\mathop{\mathrm{R}^{#1}\hspace{-2pt}}\nolimits}

\newcommand{\simar}{\mathrel{\overset{\sim}{\rightarrow}}}%同型右矢印
\newcommand{\simarr}{\mathrel{\overset{\sim}{\longrightarrow}}}%同型右矢印
\newcommand{\simra}{\mathrel{\overset{\sim}{\leftarrow}}}%同型左矢印
\newcommand{\simrra}{\mathrel{\overset{\sim}{\longleftarrow}}}%同型左矢印

\newcommand{\hocolim}{{\mathrm{hocolim}}}
\newcommand{\indlim}[1][]{\mathop{\varinjlim}\limits_{#1}}
\newcommand{\sindlim}[1][]{\smash{\mathop{\varinjlim}\limits_{#1}}\,}
\newcommand{\Pro}{\mathrm{Pro}}
\newcommand{\Ind}{\mathrm{Ind}}
\newcommand{\prolim}[1][]{\mathop{\varprojlim}\limits_{#1}}
\newcommand{\sprolim}[1][]{\smash{\mathop{\varprojlim}\limits_{#1}}\,}

\newcommand{\Sh}{\mathrm{Sh}}
\newcommand{\PSh}{\mathrm{PSh}}

\newcommand{\rmD}{\mathrm{D}}

\newcommand{\Lloc}[1][]{\mathord{\mathcal{L}^1_{\mathrm{loc},{#1}}}}
\newcommand{\ori}{\mathord{\mathrm{or}}}
\newcommand{\Db}{\mathord{\mathscr{D}b}}

\newcommand{\codim}{\mathop{\mathrm{codim}}\nolimits}



\newcommand{\gld}{\mathop{\mathrm{gld}}\nolimits}
\newcommand{\wgld}{\mathop{\mathrm{wgld}}\nolimits}


\newcommand{\tens}[1][]{\mathbin{\otimes_{\raise1.5ex\hbox to-.1em{}{#1}}}}
\newcommand{\etens}{\mathbin{\boxtimes}}
\newcommand{\ltens}[1][]{\mathbin{\overset{\mathrm{L}}\tens}_{#1}}
\newcommand{\mtens}[1][]{\mathbin{\overset{\mathrm{\mu}}\tens}_{#1}}
\newcommand{\lltens}[1][]{{\mathop{\tens}\limits^{\mathrm{L}}_{#1}}}
\newcommand{\letens}{\overset{\mathrm{L}}{\etens}}
\newcommand{\detens}{\underline{\etens}}
\newcommand{\ldetens}{\overset{\mathrm{L}}{\underline{\etens}}}
\newcommand{\dtens}[1][]{{\overset{\mathrm{L}}{\underline{\otimes}}}_{#1}}

\newcommand{\blk}{\mathord{\ \cdot\ }}
\newcommand{\mres}[2][]{{\left.{#1}\right\rvert}_{#2}}


%\newcommand{\hocolim}{{\mathrm{hocolim}}}
%\newcommand{\indlim}[1][]{\mathop{\varinjlim}\limits_{#1}}
%\newcommand{\sindlim}[1][]{\smash{\mathop{\varinjlim}\limits_{#1}}\,}
%\newcommand{\Pro}{\mathrm{Pro}}
%\newcommand{\Ind}{\mathrm{Ind}}
%\newcommand{\prolim}[1][]{\mathop{\varprojlim}\limits_{#1}}
%\newcommand{\sprolim}[1][]{\smash{\mathop{\varprojlim}\limits_{#1}}\,}
\newcommand{\proolim}[1][]{\mathop{\text{\rm``{$\varprojlim$}''}}\limits_{#1}}
\newcommand{\sproolim}[1][]{\smash{\mathop{\rm``{\varprojlim}''}\limits_{#1}}}
\newcommand{\inddlim}[1][]{\mathop{\text{\rm``{$\varinjlim$}''}}\limits_{#1}}
\newcommand{\sinddlim}[1][]{\smash{\mathop{\text{\rm``{$\varinjlim$}''}}\limits_{#1}}\,}
\newcommand{\ooplus}{\mathop{\text{\rm``{$\oplus$}''}}\limits}
\newcommand{\bbigsqcup}{\mathop{``\bigsqcup''}\limits}
\newcommand{\bsqcup}{\mathop{``\sqcup''}\limits}
\newcommand{\dsum}[1][]{\mathbin{\oplus_{#1}}}

\newcommand{\Fct}{\mathop{\mathsf{Fct}}\nolimits}

\newcommand{\pb}[1][]{\mathbin{\mathop{\times}\limits_{#1}}}
\newcommand{\dpb}[1][]{\mathbin{\mathop{\times}\limits_{#1}}}
\newcommand{\tpb}[1][]{\mathbin{\mathop{\times}\nolimits_{#1}}}





%================================================
% 自前の定理環境
%   https://mathlandscape.com/latex-amsthm/
% を参考にした
\newtheoremstyle{mystyle}%   % スタイル名
    {5pt}%                   % 上部スペース
    {5pt}%                   % 下部スペース
    {}%              % 本文フォント
    {}%                  % 1行目のインデント量
    {\bfseries}%                      % 見出しフォント
    {.}%                     % 見出し後の句読点
    {12pt}%                     % 見出し後のスペース
    {\thmname{#1}\thmnumber{ #2}\thmnote{{\hspace{2pt}\normalfont (#3)}}}% % 見出しの書式

\theoremstyle{mystyle}
\newtheorem{AXM}{公理}[section]
\newtheorem{DFN}[AXM]{定義}
\newtheorem{THM}[AXM]{定理}
\newtheorem*{THM*}{定理}
\newtheorem{PRP}[AXM]{命題}
\newtheorem{LMM}[AXM]{補題}
\newtheorem{CRL}[AXM]{系}
\newtheorem{EG}[AXM]{例}
\newtheorem*{EG*}{例}
\newtheorem{RMK}[AXM]{注意}
\newtheorem{CNV}[AXM]{約束}
\newtheorem{CMT}[AXM]{コメント}
\newtheorem*{CMT*}{コメント}
\newtheorem{NTN}[AXM]{記号}

% 定理環境ここまで
%====================================================

\usepackage{framed}
\definecolor{lightgray}{rgb}{0.75,0.75,0.75}
\renewenvironment{leftbar}{%
  \def\FrameCommand{\textcolor{lightgray}{\vrule width 4pt} \hspace{10pt}}% 
  \MakeFramed {\advance\hsize-\width \FrameRestore}}%
{\endMakeFramed}
\newenvironment{redleftbar}{%
  \def\FrameCommand{\textcolor{lightgray}{\vrule width 1pt} \hspace{10pt}}% 
  \MakeFramed {\advance\hsize-\width \FrameRestore}}%
 {\endMakeFramed}



\renewcommand{\Re}{\mathop{\mathrm{Re}}\nolimits}

% =================================





% ---------------------------
% new definition macro
% ---------------------------
% 便利なマクロ集です

% 内積のマクロ
%   例えば \inner<\pphi | \psi> のように用いる
\def\inner<#1>{\langle #1 \rangle}

% ================================
% マクロを追加する場合のスペース 

%=================================



% 位相空間
\newcommand{\Int}[1][]{\mathop{\mathrm{Int}}\nolimits_{#1}}
\newcommand{\Cl}[1][]{\mathop{\mathrm{Cl}}\nolimits_{#1}}
\newcommand{\Bd}[1][]{\mathop{\mathrm{Bd}}\nolimits_{#1}}
\newcommand{\bd}{\mathord{\partial}}
\newcommand{\Fr}[1][]{\mathop{\mathrm{Fr}}\nolimits_{#1}}


\newcommand{\mtan}[1][]{T\hspace{-1.75pt}{#1}}
\newcommand{\mtp}[1][]{{}^{t\!}{#1}} % transpose of the morphism
%\newcommand{\mpb}[1][]{{}^{t\!}{#1}'} % pull-back of the morphism
\newcommand{\mpb}[1][]{{#1}_{\pi}} % pull-back of the morphism
\newcommand{\mdiff}[1][]{{#1}_{d}} % pull-back of the morphism
\newcommand{\mptan}[1][]{{#1}_{\tau}} % pull-back of the morphism


%\newcommand{\pb}[1][]{\mathbin{\mathop{\times}\limits_{#1}}}

\newcommand{\cotb}[1][]{T^{\ast}{#1}}
\newcommand{\conm}[2][]{T_{#1}^{\hspace{0.7pt}\ast}{#2}}
\newcommand{\norb}[2][]{T_{#1}{#2}}

\newcommand{\cotz}[1][]{T^{\ast}_{#1}{#1}}

\newcommand{\cotn}[1][]{\dot{T}^{\ast}{#1}}


\newcommand{\muhom}{\mu hom}
\newcommand{\muHom}[1][]{\mathrm{Hom}^\mu_{\raise1.5ex\hbox to.1em{}#1}}

\newcommand{\mucat}[2][]{\mathsf{D}^{\mathrm{b}}_{#1}(#2)}
%\newcommand{\mucat}[2][]{\mathsf{D}^{\mathrm{b}}_{#1}(#2)}

\newcommand{\conv}[1][]{\mathbin{\mathop{\circ}\limits_{#1}}}


\newcommand{\lopen}{\mathopen{]}}
\newcommand{\lclosed}{\mathopen{[}}
\newcommand{\ropen}{\mathclose{[}}
\newcommand{\rclosed}{\mathopen{]}}

\newcommand{\sm}{\raisebox{2.33pt}{~\rule{6.4pt}{1.3pt}~}}
\newcommand{\smi}{\smallsetminus}

\newcommand{\smid}{\mathrel{;}}


% ----------------------------
% documenet 
% ----------------------------
% 以下, 本文の執筆スペースです. 
% Your main code must be written between 
% begin document and end document.
% ---------------------------

\title{Truncated Microlocal Morse Theory}
\author{Toshihiro Oshiba}
\date{\today\footnote{First draft: 2026/03/29}}
\begin{document}
\maketitle
\section{Truncated microsupport}
\subsection*{Functorial Properties} 
\begin{Proposition+}
  Let \(f\colon X\to Y\) be a morphism of real analytic manifolds and let \(F\in\Dompb(\kk_X)\) such that \(f\) is proper on \(\supp(F)\). 
  Then for any \(k\in\zz\), 
  \begin{equation}\label{eq:t-push}
    \muS[k](\Rder{f_{\ast}}F)\subset
    \mpb[f]\mdiff[f]^{-1}(\muS[k](F)).
  \end{equation}
  The equality holds in case \(f\) is a closed embedding.
\end{Proposition+}

\section{Truncated microlocal Morse theory}

Let \(f\colon Y\to X\) be a morphism of manifolds, and 
\(\varphi\) a real valued function of class \(C^1\). 
For a real number \(t\in \rr\), we set 
\begin{align*}
  Y_t
  &\coloneqq \left\{y\in Y\smid \varphi(y)< t\right\} 
  = \varphi^{-1}(\lopen-\infty,t\ropen),\\
  \widetilde{Y_t}
  &\coloneqq \left\{y\in Y\smid \varphi(y)\leqq t\right\} 
  = \varphi^{-1}(\lopen-\infty,t\rclosed).
\end{align*}
Let us define the following morphisms:
\[\xymatrix{
  Y_t\ar[r]^-{j_t}
  \ar[rd]_-{f_t}
  &
  Y
  \ar[d]^-{f}\\
  &X,
}\quad\xymatrix{
  \widetilde{Y_t}\ar[r]^-{\tilde{j}_t}
  \ar[rd]_-{f_t}
  &
  Y
  \ar[d]^-{f}\\
  &X.
}\]

\begin{Proposition+}
  Let \(G\in\Dompb(Y)\), and assume \(
    f_{t}\colon \supp(G)\cap \overline{Y_t}\to X
  \) is proper for all \(t\in\rr\). 
  Let \(k\) be an integer and \(t_\circ\) a real number.
  \begin{enumerate}[(i)]
    \item Assume that, 
    for any \(y\in Y\smi \widetilde{Y}_{t_{\circ}}\), 
    we have \[
      d\varphi(y)\notin \left(
        \muS[k](G)+ \mdiff[f]\left(
          Y\pb[X]\cotb[X]
        \right)
      \right).
    \]
    Then the following statements hold.
    \begin{enumerate}[(a)]
        \item For \(t>t_\circ\), the natural morphism
        \[
          \Rder[j]{f_{\ast}}G \to \Rder[j]{{f_{t}}_{\ast}}j^{-1}_{t}G
        \]
        is isomorphic for \(j<k\) and injective for \(j=k\). 
        \item[(a\('\))] For \(t\geqq t_\circ\), the natural morphism
        \[
          \Rder[j]{\tilde{f}_{\ast}}G 
          \to 
          \Rder[j]\tilde{f_{t}}_{\ast}\tilde{j}^{-1}_{t}G
        \]
        is isomorphic for \(j<k\) and injective for \(j=k\). 
        \item We have \begin{align*}
          \muS[k](\Rder[]{f_{\ast}}G)
          &\subset
          \mpb[f]\left(\mdiff[f]^{-1}\left(
            \muS[k](G)\cap \pi^{-1}(\widetilde{Y}_{t_{\circ}})
          \right)\right)\\
          &=\mpb[f]\mdiff[f]^{-1}\left(
          \muS[k](G)
          \right).
        \end{align*}
    \end{enumerate}
    \item Assume that, 
    for any \(y\in Y\smi \widetilde{Y}_{t_{\circ}}\), 
    we have \[
      -d\varphi(y)\notin \left(
        \muS[k](G)+ \mdiff[f]\left(
          Y\pb[X]\cotb[X]
        \right)
      \right).
    \]
    Then \begin{align}
      \tag{c}
      \Rder[j]{f_{!}}G
      &\cong
      \Rder[j]{{f_{t}}_{!}}j^{-1}_{t}G
      \quad
      \text{ for \(j<k\) and \(t>t_\circ\)},\\
      \tag{c\(_k\)}
      \Rder[k]{f_{!}}G
      &\hookrightarrow
      \Rder[k]{{f_{t}}_{!}}j^{-1}_{t}G
      \quad
      \text{ for \(t>t_\circ\)},\\
      \tag{c\('\)}
      \Rder[j]{\tilde{f}_{!}}G
      &\cong
      \Rder[j]\tilde{f_{t}}_{\ast}\tilde{j}^{!}_{t}G
      \quad
      \text{ for \(j<k\) and \(t\geqq t_\circ\)},\\
      \tag{c\(_k'\)}
      \Rder[k]{\tilde{f}_{!}}G
      &\hookrightarrow
      \Rder[k]\tilde{f_{t}}_{\ast}\tilde{j}^{!}_{t}G
      \quad
      \text{ for \(t\geqq t_\circ\)},\\
      \tag{d}
      \muS[k](\Rder[]{f_{!}}G)
      &\subset
      \mpb[f]\left(\mdiff[f]^{-1}\left(
        \muS[k](G)\cap \pi^{-1}(\widetilde{Y}_{t_{\circ}})
      \right)\right)\\
      %&
      \notag&=\mpb[f]\mdiff[f]^{-1}\left(
        \muS[k](G)
      \right).
    \end{align}
  \end{enumerate}
\end{Proposition+}

\begin{proof}
  Let \(\tilde{f}\) be the morphism \(
    (f,\varphi)\colon Y\to X\times \rr
  \), and let \(q\) and \(\tilde{\varphi}\) be 
  the first and second projections 
  defined on \(X\times \rr\) respectively. 
  Then \(\tilde{f}\) is proper on \(\supp(G)\) and 
  \(f=q\circ \tilde{f}\). 
  Let \(i_t\) be the open embedding \(
    X\times\lopen -\infty, t\ropen\hookrightarrow X\times \rr
  \) and \(\tilde{i}_t\) the closed embedding \(
    X\times\lopen -\infty, t\rclosed\hookrightarrow X\times \rr
  \). Set \(q_t=q\circ i_t\) and \(\tilde{q}_t=q\circ\tilde{i}_t\). 
  These morphisms can be visualized 
  by the following diagrams, in both of which the left squares are cartesian.
  \[\xymatrix@C=30pt@R=30pt{
    Y
      \ar[r]^-{\tilde{f}}
      \ar@{}[rd]|\square
    &
    X\times\rr
      \ar[r]^-{q}
    &
    X\\
    Y_{t}
      \ar@{^(->}[u]^-{j_t}
      \ar[r]^-{\hat{f}_{t}}
    &
    X\times \lopen-\infty,t\ropen
      \ar@{^(->}[u]_-{i_{t}}
      \ar[ur]_-{q_{t}}
    &
  },\quad
  \xymatrix@C=30pt@R=30pt{
    Y
      \ar[r]^-{\tilde{f}}
      \ar@{}[rd]|\square
    &
    X\times\rr
      \ar[r]^-{q}
    &
    X\\
    \widetilde{Y}_{t}
      \ar@{^(->}[u]^-{\tilde{j}_t}
      \ar[r]^-{\check{f}_{t}}
    &
    X\times \lopen-\infty,t\rclosed
      \ar@{^(->}[u]_-{\tilde{i}_{t}}
      \ar[ur]_-{\tilde{q}_{t}}
    &
  }.\]
  Then, setting \(\widetilde{G}\coloneqq \Rder[]{f_{\ast}}G\), 
  we have \begin{align*}
    \Rder[]{f_{\ast}}G
    &\cong 
    \Rder[]{q_{\ast}}\widetilde{G},
    &
    \Rder[]{\tilde{f}_{\ast}}G
    &\cong 
    \Rder[]{\tilde{q}_{\ast}}\widetilde{G},
    \\
    \Rder[]{{f_{t}}_{\ast}}j_{t}^{-1}G
    &\cong 
    \Rder[]{{q_{t}}_{\ast}}i_{t}^{-1}\widetilde{G},
    &
    \Rder[]{\tilde{f_{t}}_{\ast}}\tilde{j}_{t}^{-1}G
    &\cong 
    \Rder[]{\tilde{q_{t}}_{\ast}}\tilde{i}_{t}^{-1}\widetilde{G},
    \\
    \Rder[]{{f_{t}}_{!}}j_{t}^{-1}G
    &\cong 
    \Rder[]{{q_{t}}_{!}}i_{t}^{-1}\widetilde{G},
    &
    \Rder[]{\tilde{f_{t}}_{\ast}}\tilde{j}_{t}^{!}G
    &\cong 
    \Rder[]{\tilde{q_{t}}_{\ast}}\tilde{i}_{t}^{!}\widetilde{G}.
  \end{align*}
  Thus taking \(j\)-th cohomology, we have\begin{align*}
    \Rder[j]{f_{\ast}}G
    &\cong 
    \Rder[j]{q_{\ast}}\widetilde{G},
    &
    \Rder[j]{\tilde{f}_{\ast}}G
    &\cong 
    \Rder[j]{\tilde{q}_{\ast}}\widetilde{G},
    \\
    \Rder[j]{{f_{t}}_{\ast}}j_{t}^{-1}G
    &\cong 
    \Rder[j]{{q_{t}}_{\ast}}i_{t}^{-1}\widetilde{G},
    &
    \Rder[j]{\tilde{f_{t}}_{\ast}}\tilde{j}_{t}^{-1}G
    &\cong 
    \Rder[j]{\tilde{q_{t}}_{\ast}}\tilde{i}_{t}^{-1}\widetilde{G},
    \\
    \Rder[j]{{f_{t}}_{!}}j_{t}^{-1}G
    &\cong 
    \Rder[j]{{q_{t}}_{!}}i_{t}^{-1}\widetilde{G},
    &
    \Rder[j]{\tilde{f_{t}}_{\ast}}\tilde{j}_{t}^{!}G
    &\cong 
    \Rder[j]{\tilde{q_{t}}_{\ast}}\tilde{i}_{t}^{!}\widetilde{G}.    
  \end{align*}
  Therefore, it is enough to prove the result for \(\widetilde{G}\), \(q\), \(i_t\), \(\tilde{\varphi}\) instead of \(G\), \(f\), \(j_t\), \(\varphi\). 

  (i) 
  It follows from the hypothesis that
  \begin{equation}\label{eq:tilde-estimate}
    \muS[k](\widetilde{G})\cap\left(
      \cotb[X]\times\left\{(t;\tau);t>t_{\circ}\right\}
    \right)
    \subset
    \cotb[X]\times\left\{(t;\tau);\tau\leqq0\right\}.
  \end{equation}
  Indeed, by hypothesis, 
  for any \((y;\eta)\in\muS[k](G)\) 
  and \((f(x);\xi)\in Y\tpb[X]\cotb[X]\), 
  we have \begin{equation}\label{eq:pb-conic}
    d\varphi(y)\neq \eta + f^{\ast}\xi.
  \end{equation}
  If \((x,t;\xi,\tau)\) is a point of \(
    \muS[k](\widetilde{G})\cap\left(
      \cotb[X]\times\left\{(t;\tau);t>t_{\circ}\right\}
  \right)\), 
  then, by the formula \eqref{eq:t-push}, 
  there exists \((y;\eta)\in\muS[k](G)\) such that \[
    \tilde{f}(y)=(x,t) 
    \text{ and } 
    \eta=\tilde{f}^{\ast}(\xi,\tau).
  \] For this \(\eta\), we have \[
    \eta
    =\tilde{f}^{\ast}(\xi,\tau)
    =f^{\ast}\xi +\varphi^{\ast}\tau
    =f^{\ast}\xi +\tau d\varphi(y).
  \]
  If \(\tau>0\), we get \[
    d\varphi(y)=\frac{1}{\tau}\eta-\frac{1}{\tau}f^{\ast}\xi
  \] and this contradicts \eqref{eq:pb-conic} 
  because \(\muS[k](G)\) is conic. 

  By \eqref{eq:tilde-estimate}, 
  Let \(U\) be an open subset of \(X\) contained 
  in a local chart, and convex in such a chart. 

\end{proof}













%\section{a}
%===============================================
% 参考文献スペース
%===============================================
\begin{thebibliography}{20} 
  \bibitem[GKS12]{GKS12} St\'ephane Guillermou, 
  Masaki Kashiwara, Pierre Schapira, 
  \textit{Sheaf Quantization of Hamiltonian Isotopies and Applications to Nondisplaceability Problems}, 
  Duke. Math. J. 161, No. 2, pp.201--245, 2012.  
  \bibitem[KFS03]{KFS03} Masaki Kashiwara, Teresa Monteiro Fernandes, Pierre Schapira, 
    \textit{Truncated Microsupport and Holomorphic Solutions 
    of \(D\)-modules}, 
    Ann. Scient. \'Ec. Norm. Sup., s\'erie 4, tome 36, 
    pp.583--399, 2003.
  \bibitem[KS90]{KS90} Masaki Kashiwara, Pierre Schapira, 
    \textit{Sheaves on Manifolds}, 
    Grundlehren der Mathematischen Wissenschaften, 292, Springer, 1990.
  \bibitem[ST92]{ST92} Pierre Schapira, Nobuyuki Tose, 
    \textit{Morse Inequalities for \(\rr\)-constructible Sheaves}, 
    Adv. Math. 93, pp.1--8, 1992. 
\end{thebibliography}
%===============================================


\end{document}
